Within the field of a galaxy cluster at a redshift of $z = 0.500$, a galaxy
which looks like a normal elliptical is observed, with an apparent magnitude in
the $B$ filter $m_B = 20.40$\,mag.

The luminosity distance corresponding to a redshift of $z = 0.500$ is
$d_L = 2754$\,Mpc.

The spectral energy distribution (SED) of elliptical galaxies in the wavelength
range 250\,nm to 500\,nm is adequately approximated by the formula:
\[ L_\lambda(\lambda) \propto \lambda^4 \]
(i.e., the spectral density of the object's luminosity, known also as the
monochromatic luminosity, is proportional to $\lambda^4$.)

\begin{description}
\item[(a)] What is the absolute magnitude of this galaxy in the $B$ filter?
\item[(b)] Can it be a member of this cluster? (write YES or NO alongside your
  final calculation)
\end{description}

Hints: Try to establish a relation that describe the dependence of the spectral
density of flux on distance for small wavelength interval. Normal elliptical
galaxies have maximum absolute magnitude equal to $-22$\,mag.
